% Generated by utilities/generate_user_guide.py. Edit routine wording in utilities/descriptions.yaml where possible. This is a content-only LaTeX file; include it from a top-level document. Long-form appendices are sourced from docs/appendices/; references are maintained in docs/references.bib.

\section{Introduction}

Establishing an in vitro-in vivo correlation (IVIVC) can support the evaluation of medical devices that include absorbable polymers by relating in vivo changes in material properties over time to corresponding in vitro test results, such as degradation profiles. IVIVCs may help evaluate manufacturing or material changes, support predictions of in vivo performance when modifying an existing device design, and reduce reliance on animal studies during preclinical development. This tool is intended to facilitate IVIVC development. The user supplies paired in vitro and in vivo fitting data, and may optionally supply an in vitro prediction dataset. The app applies selected preprocessing steps, fits and compares multiple correlation approaches and model families, evaluates selected performance metrics and cross-validation results, generates fitting and prediction plots, and produces report-ready summaries that can be reviewed and incorporated into external documentation.

For general background on credibility assessment and reporting of computational modeling studies in medical device submissions, see FDA guidance on computational modeling credibility and reporting \cite{fda_modeling_guidance,fda_reporting_guidance}.

\section{What the app does}

\begin{itemize}
\item Applies selected normalization, scaling, and interpolation options to prepare fitting and prediction data.
\item Fits user-selected IVIVC approaches and model families to paired in vitro and in vivo data.
\item Calculates selected goodness-of-fit, model-comparison, and cross-validation metrics.
\item Generates fitting, comparison, and prediction plots with uncertainty displays where supported.
\item Creates summaries for user-selected final approaches or models.
\end{itemize}

\section{Inputs}

The IVIVC app requires paired in vitro and in vivo degradation data, along with user-selected analysis settings that define how the data are processed, modeled, and summarized. At a high level, the inputs include:

\begin{itemize}
\item a spreadsheet containing in vitro and in vivo time-course degradation data for fitting correlations;
\item optionally, a spreadsheet containing in vitro data from which to generate predictions using fitted correlations;
\item selection of one or more IVIVC modeling approaches;
\item selection of candidate models to evaluate, where applicable;
\item optionally, preprocessing options, such as normalization, interpolation, or scaling choices;
\item selected performance metrics to display in the results and reports; by default, the app selects \(R^2\), RMSE, and AIC\textsubscript{c} as a simple starting set that has been useful in internal testing with real datasets but may not be optimal for every analysis
\item optionally, cross-validation settings for applicable parametric models; and
\item grid-search settings for model initialization.
\end{itemize}

The sections below provide additional detail on input file structure, required data columns, preprocessing options, model-selection inputs, and optional analysis settings.

\section{Outputs}

The IVIVC app generates interactive results and downloadable reports that summarize fitted model behavior, model-comparison metrics, diagnostic information, and user-selected analysis settings. Output availability depends on the selected approaches, whether fitting succeeds, and whether a prediction dataset was supplied.

\begin{longtable}{@{}>{\raggedright\arraybackslash}p{0.24\textwidth}>{\raggedright\arraybackslash}p{0.68\textwidth}@{}}
\toprule
\textbf{Output} & \textbf{Description} \\
\midrule
\endfirsthead
\toprule
\textbf{Output} & \textbf{Description} \\
\midrule
\endhead
Preprocessing outputs & Processed data tables and plots showing the data used for fitting and prediction. \\
Model fitting outputs & Parameter estimates, parameter standard errors where available, fitted plots, residual (error) plots, and selected performance metrics. \\
Model comparison outputs & Tables comparing selected approaches, models, metrics, and cross-validation summaries. \\
Prediction outputs & Prediction plots and downloadable tables when a prediction dataset is supplied and the selected method can be applied. \\
Selected model \& summary of outputs & User-selected report summaries and optional Word report export for incorporation into external documentation. \\
\bottomrule
\end{longtable}

\section{Input datasets}

\begin{longtable}{@{}>{\raggedright\arraybackslash}p{0.24\textwidth}>{\raggedright\arraybackslash}p{0.68\textwidth}@{}}
\toprule
\textbf{Input} & \textbf{Description} \\
\midrule
\endfirsthead
\toprule
\textbf{Input} & \textbf{Description} \\
\midrule
\endhead
Fitting dataset & Required four-column dataset used to establish the in vitro/in vivo relationship. Columns are interpreted as in vitro time, in vitro value, in vivo time, and in vivo value. \\
Prediction dataset & Optional two-column in vitro dataset used after fitting to generate in vivo predictions. Columns are interpreted as in vitro time and in vitro value. \\
\bottomrule
\end{longtable}

For Excel workbooks, select the sheet containing the relevant data. For CSV files, the first columns are read directly in the order described above.

\section{Analysis approaches and models}

\begin{longtable}{@{}>{\raggedright\arraybackslash}p{0.18\textwidth}>{\raggedright\arraybackslash}p{0.46\textwidth}>{\raggedright\arraybackslash}p{0.28\textwidth}@{}}
\toprule
\textbf{Approach} & \textbf{Description} & \textbf{Available models} \\
\midrule
\endfirsthead
\toprule
\textbf{Approach} & \textbf{Description} & \textbf{Available models} \\
\midrule
\endhead
Approach 1: Direct Correlation & Approach 1 fits a direct relationship between matched in vitro and in vivo time points after preprocessing. The fitted model is then used as a time-scaling function that maps an in vitro time to a predicted in vivo time. & Double Exponential, Exponential, Linear, Linear (zero intercept), Power \\
Approach 2: Independent Model Fitting & Approach 2 fits the same model family independently to the in vitro and in vivo datasets. With fitting data alone, the report summarizes the two fitted models and their fitted-curve uncertainty. When prediction data are supplied, predictions can be generated by comparing the fitted in vivo and in vitro responses or by comparing fitted characteristic time constants when those values are available. & Double Exponential, Exponential, Mass Loss (Gompertz), Mass Loss (One-Phase), Mass Loss (Surface Erosion), Mass Loss (Two-Phase), Power, Stretched Exponential \\
Approach 3: Direct Mapping & Approach 3 uses direct interpolation-based mappings between datasets rather than fitting a parametric model. It can describe value-ratio behavior over time and time-ratio behavior over value ranges. When prediction data are supplied, those mappings can be used to scale prediction values or prediction times. & No parametric model selection \\
\bottomrule
\end{longtable}

\subsection{Candidate model functional forms}

\begin{longtable}{@{}>{\raggedright\arraybackslash}p{0.16\textwidth}>{\raggedright\arraybackslash}p{0.14\textwidth}>{\raggedright\arraybackslash}p{0.36\textwidth}>{\raggedright\arraybackslash}p{0.26\textwidth}@{}}
\toprule
\textbf{Model} & \textbf{Used in} & \textbf{Functional form} & \textbf{Notes} \\
\midrule
\endfirsthead
\toprule
\textbf{Model} & \textbf{Used in} & \textbf{Functional form} & \textbf{Notes} \\
\midrule
\endhead
Double Exponential & Approach 1, Approach 2 & \(y = a e^{b x} + c e^{d x}\) & Sum of two exponential terms. It can represent profiles with two apparent rates, but the added flexibility may require more data to estimate stably. \\
Exponential & Approach 1, Approach 2 & \(y = a e^{b x}\) & Two-parameter exponential relationship. It can describe monotonic change that accelerates or decelerates exponentially over the fitted range. \\
Linear & Approach 1 & \(y = a + b x\) & Straight-line empirical relationship with an intercept and slope. It is appropriate when the in vitro/in vivo relationship is approximately linear over the fitted range. \\
Linear (zero intercept) & Approach 1 & \(y = a x\) & Straight-line empirical relationship constrained to pass through zero. It may be useful when a zero in vitro value should correspond to a zero in vivo value under the selected preprocessing and response scale. \\
Mass Loss (Gompertz) & Approach 2 & \(y = y_0 - A \exp\{-\exp[(e r_{max}/A)(t_{lag}-x)+1]\}\) & Mass-remaining model based on the modified Gompertz curve. It is intended for asymmetric sigmoidal mass-loss profiles with an apparent lag or plateau, a maximum loss-rate region, and a terminal loss extent. \\
Mass Loss (One-Phase) & Approach 2 & \(y = y_0 - A/(1 + e^{-(x-t_{50})/s})\) & Mass-remaining model with a delayed sigmoidal loss phase. It is intended for mass-loss profiles that begin constant and show a later drop. \\
Mass Loss (Surface Erosion) & Approach 2 & \(y = y_0 \max(1-kx,0)^n\) & Mass-remaining model for surface-controlled erosion. The fitted exponent n allows the profile shape to adapt to slab-like, cylindrical, spherical, or empirical geometry-like erosion behavior. \\
Mass Loss (Two-Phase) & Approach 2 & \(y = y_0 - B(1-e^{-k_b x}) - D/(1+e^{-(x-t_{50})/s})\) & Mass-remaining model with an optional early burst-loss phase followed by a delayed sigmoidal loss phase. It is intended for mass-loss profiles that may show an initial drop, an intermediate plateau, and a later drop. \\
Power & Approach 1, Approach 2 & \(y = a x^{b}\) & Power-law empirical relationship. It can describe nonlinear monotonic relationships where the response changes as a fitted power of time or the mapped predictor. \\
Stretched Exponential & Approach 2 & \(y = a e^{b x^{c}}\) & Exponential relationship with a fitted stretching exponent. It can describe non-single-rate behavior, but the additional shape parameter can increase sensitivity when data are sparse. \\
\bottomrule
\end{longtable}

\textbf{Mass-loss model source notes}

The surface-erosion mass-loss form is based on geometry-controlled surface erosion concepts such as the erodible slab, cylinder, and sphere model \cite{hopfenberg1976}. The one-phase mass-loss model uses a logistic/sigmoidal transition to represent a delayed loss phase \cite{tsoularis2002}. The two-phase mass-loss model combines an early burst-like loss term with a delayed sigmoidal loss phase. Burst and delayed/multiphase behavior in degradable polymer systems is discussed in the polymer erosion and PLGA release literature \cite{gopferich1997,fredenberg2011,allison2008}. The Gompertz mass-loss model uses a modified Gompertz-style sigmoidal form with lag and maximum-rate parameters \cite{zwietering1990}.

\section{Preprocessing options}

\subsection{Normalization}

\begin{longtable}{@{}>{\raggedright\arraybackslash}p{0.22\textwidth}>{\raggedright\arraybackslash}p{0.20\textwidth}>{\raggedright\arraybackslash}p{0.50\textwidth}@{}}
\toprule
\textbf{Option} & \textbf{Internal key} & \textbf{Description} \\
\midrule
\endfirsthead
\toprule
\textbf{Option} & \textbf{Internal key} & \textbf{Description} \\
\midrule
\endhead
Min-max normalization & min\_max & Scales response values to the 0-1 range using the observed minimum and maximum within each dataset. \\
Initial-value normalization & initial\_value & Divides response values by the first response value in the corresponding dataset. \\
\bottomrule
\end{longtable}

\subsection{Scaling}

\begin{longtable}{@{}>{\raggedright\arraybackslash}p{0.22\textwidth}>{\raggedright\arraybackslash}p{0.20\textwidth}>{\raggedright\arraybackslash}p{0.50\textwidth}@{}}
\toprule
\textbf{Option} & \textbf{Internal key} & \textbf{Description} \\
\midrule
\endfirsthead
\toprule
\textbf{Option} & \textbf{Internal key} & \textbf{Description} \\
\midrule
\endhead
Log10 time scaling & log\_x & Applies a base-10 logarithm to time values before fitting or interpolation. \\
Log10 response scaling & log\_y & Applies a base-10 logarithm to response values before fitting or interpolation. \\
\bottomrule
\end{longtable}

\subsection{Interpolation}

\begin{longtable}{@{}>{\raggedright\arraybackslash}p{0.22\textwidth}>{\raggedright\arraybackslash}p{0.20\textwidth}>{\raggedright\arraybackslash}p{0.50\textwidth}@{}}
\toprule
\textbf{Option} & \textbf{Internal key} & \textbf{Description} \\
\midrule
\endfirsthead
\toprule
\textbf{Option} & \textbf{Internal key} & \textbf{Description} \\
\midrule
\endhead
No interpolation & none & Uses only exact raw alignments. No values are estimated between observations; unavailable aligned positions are marked as missing. \\
Shared-time interpolation & default & Interpolates both datasets onto the common set of experimental time points, using linear interpolation and marking out-of-range values as unavailable. \\
Densified interpolation & alternative & Creates additional linearly spaced points between adjacent observations before interpolation. This can smooth displays and make fitting more stable, but adds more interpolation-derived points that skew statistical comparisons. \\
\bottomrule
\end{longtable}

\section{Relative model evidence}

Relative model evidence compares candidate parametric models fitted to the same dataset under the same preprocessing settings and IVIVC approach. The evidence table is based only on the information-criterion metrics selected by the user: AIC, AIC\textsubscript{c}, and/or BIC. These criteria are relative; their absolute values are not interpreted directly, and negative values are allowed. Table definitions: n is the number of fitted observations; k is the number of fitted model parameters; Delta is the difference between a model's criterion value and the lowest value in the displayed candidate set; weight is the normalized relative support for that model under the criterion; and evidence ratio compares the highest weight in the set with the model's weight. Delta values near 0 indicate similar relative support; larger delta values indicate lower support relative to the leading candidate under that criterion. Similar relative evidence does not mean the models are scientifically equivalent. Final model selection should also consider prediction error, cross-validation behavior, residual patterns, parameter uncertainty, model simplicity, and biological or degradation-mechanism plausibility. The app reports these summaries to support user review. It does not select, recommend, or validate a model.

For additional background on information criteria, model-selection evidence, and evidence-weight interpretation, see references \cite{akaike1974,schwarz1978,burnham2002,kass1995,wagenmakers2004}.

\begin{longtable}{@{}>{\raggedright\arraybackslash}p{0.24\textwidth}>{\raggedright\arraybackslash}p{0.68\textwidth}@{}}
\toprule
\textbf{Evidence summary} & \textbf{Interpretation} \\
\midrule
\endfirsthead
\toprule
\textbf{Evidence summary} & \textbf{Interpretation} \\
\midrule
\endhead
Delta AIC\textsubscript{c} / Delta AIC / Delta BIC & Difference from the lowest value in the displayed candidate set. Values near 0 indicate similar relative support; larger values indicate lower support relative to the leading candidate under that criterion. \\
AIC\textsubscript{c} or BIC weight & Normalized relative-support value across the candidate models in the displayed evidence table. Weights are relative, not absolute probabilities that a model is true. \\
Evidence ratio & Ratio of the highest evidence weight to the model's evidence weight. Larger ratios indicate less relative support compared with the highest-weighted model in the set. \\
Interpretation badge & Neutral rule-of-thumb label based on the displayed delta values. The label supports user review and is not an app recommendation. \\
\bottomrule
\end{longtable}

\section{Cross-validation options}

Cross-validation provides an approximate check of how a fitted model performs on held-out data. It is most useful when there are enough observations to fit the model repeatedly while leaving some observations out for evaluation.

The available schemes are Leave-one-timepoint-out, Leave-contiguous-block-out, Shuffle split, K-fold, and None. Leave-one-timepoint-out holds out one observation at a time, Leave-contiguous-block-out holds out adjacent time points, Shuffle split uses repeated random holdout sets, K-fold uses ordered fold-based holdouts, and None disables cross-validation. The block-size setting for Leave-contiguous-block-out controls how many consecutive observations are held out in each split. Larger blocks may leave too few observations for stable fitting, especially in small datasets. For small datasets, cross-validation summaries should be interpreted as exploratory diagnostics rather than definitive validation results.

Leave-one-timepoint-out is the default cross-validation scheme. It is a practical starting point for small time-course datasets because each split uses all but one observation for fitting. Leave-contiguous-block-out and K-fold can be used as more demanding time-region holdout checks, while Shuffle split is retained as a familiar exploratory option.

For background on cross-validation procedures and their use in model assessment and selection, particularly for for time series, see references \cite{arlot2010, bergmeir2012, roberts2017, bergmeir2018}.

\begin{longtable}{@{}>{\raggedright\arraybackslash}p{0.20\textwidth}>{\raggedright\arraybackslash}p{0.20\textwidth}>{\raggedright\arraybackslash}p{0.52\textwidth}@{}}
\toprule
\textbf{Scheme} & \textbf{Internal key} & \textbf{Description} \\
\midrule
\endfirsthead
\toprule
\textbf{Scheme} & \textbf{Internal key} & \textbf{Description} \\
\midrule
\endhead
Leave-one-timepoint-out & leave\_one\_timepoint\_out & Holds out one time point at a time and fits the model using all remaining time points. This is the default because it uses as much data as possible for fitting while still providing a held-out prediction diagnostic for small time-course datasets. \\
Leave-contiguous-block-out & leave\_contiguous\_block\_out & Holds out adjacent time points as a block in x-axis order. The block-size setting controls how many adjacent observations are held out in each split. The maximum useful block size depends on the number of observations. Larger blocks may leave too few observations for stable fitting, especially in small datasets. \\
Shuffle split & shuffle\_split & Repeatedly creates random training and holdout subsets. This can be useful as an exploratory sensitivity check when the dataset has enough observations to support repeated resampling. The splits, test-size, and random-seed settings control the repeated random holdouts. \\
K-fold & kfold & Divides the ordered dataset into a specified number of folds, fits the model on all but one fold, and evaluates on the held-out fold. With small k, this can act as a more demanding time-region holdout test. \\
None & none & Disables cross-validation. This may be appropriate when the dataset is too small for meaningful holdout validation or when only fitted-model summaries are needed. \\
\bottomrule
\end{longtable}

\begin{longtable}{@{}>{\raggedright\arraybackslash}p{0.22\textwidth}>{\raggedright\arraybackslash}p{0.22\textwidth}>{\raggedright\arraybackslash}p{0.48\textwidth}@{}}
\toprule
\textbf{Parameter} & \textbf{Internal key} & \textbf{Description} \\
\midrule
\endfirsthead
\toprule
\textbf{Parameter} & \textbf{Internal key} & \textbf{Description} \\
\midrule
\endhead
Cross-validation scheme & cv\_scheme & Selects the splitting method used to estimate held-out performance for fitted parametric models. Leave-one-timepoint-out is the default. \\
CV splits/folds/block size & cv\_n\_splits & For Shuffle split, the number of repeated random train/test splits. For K-fold, the number of folds. For Leave-contiguous-block-out, the number of adjacent time points held out per split. The setup page does not enforce a maximum block size because the uploaded data are not read until Analyze is selected; larger blocks may leave too few observations for stable fitting. This setting is not used for Leave-one-timepoint-out or None. \\
Shuffle split test size & cv\_test\_size & The fraction of observations held out for testing in each Shuffle split repeat. Smaller values leave more observations for fitting; larger values make the holdout test more demanding but leave fewer observations for training. \\
CV random seed & cv\_random\_state & Seed used to make randomized cross-validation splits reproducible. This applies to Shuffle split. \\
\bottomrule
\end{longtable}

\section{Grid-search initialization}

The app conducts an automated grid search to identify suitable starting parameter values for each model before nonlinear least-squares fitting, reducing the need for manual parameter tuning. Candidate values are sampled within the configured parameter bounds and are used only to initialize the final curve fit.

This initialization step improves convergence and reduces sensitivity to poor starting values. It is not itself the final fitting algorithm and does not guarantee that the final fitted parameters are the global optimum.

\begin{longtable}{@{}>{\raggedright\arraybackslash}p{0.24\textwidth}>{\raggedright\arraybackslash}p{0.24\textwidth}>{\raggedright\arraybackslash}p{0.44\textwidth}@{}}
\toprule
\textbf{Parameter} & \textbf{Internal key} & \textbf{Description} \\
\midrule
\endfirsthead
\toprule
\textbf{Parameter} & \textbf{Internal key} & \textbf{Description} \\
\midrule
\endhead
Grid-search random starts & grid\_search\_num\_points & Number of candidate starting parameter sets evaluated before nonlinear least-squares fitting. Higher values may improve robustness but increase runtime. \\
Grid-search CPU cores & grid\_search\_num\_cores & Number of CPU cores used for the starting-point search. Blank or Auto uses the app default. \\
Grid-search parameter minimum & grid\_search\_param\_min & Lower bound for randomly sampled candidate starting parameter values. \\
Grid-search parameter maximum & grid\_search\_param\_max & Upper bound for randomly sampled candidate starting parameter values. \\
Grid-search random seed & grid\_search\_random\_state & Seed used to make randomized starting-point sampling reproducible. Changing the seed changes the sampled candidates while keeping the same bounds and number of starts. \\
\bottomrule
\end{longtable}

\section{Performance metrics}

By default, the app selects \(R^2\), RMSE, and AIC\textsubscript{c}. This compact default set has been simple and useful in internal testing with real datasets, but it is not necessarily the best set for every dataset or intended use. Users may select additional metrics when they provide useful context for a specific analysis.

\begin{longtable}{@{}>{\raggedright\arraybackslash}p{0.17\textwidth}>{\raggedright\arraybackslash}p{0.18\textwidth}>{\raggedright\arraybackslash}p{0.20\textwidth}>{\raggedright\arraybackslash}p{0.37\textwidth}@{}}
\toprule
\textbf{Metric} & \textbf{Internal key} & \textbf{Preferred direction} & \textbf{Description} \\
\midrule
\endfirsthead
\toprule
\textbf{Metric} & \textbf{Internal key} & \textbf{Preferred direction} & \textbf{Description} \\
\midrule
\endhead
\(R^2\) & r\_squared & Higher is better & Coefficient of determination. Higher values indicate that the model explains more variance in the observed response. \\
Adjusted \(R^2\) & adjusted\_r\_squared & Higher is better & \(R^2\) adjusted for the number of fitted parameters, which can penalize unnecessary model complexity. \\
MSE & mse & Lower is better & Mean squared error between observed and fitted values. Lower values indicate smaller average squared residuals. \\
RMSE & rmse & Lower is better & Root mean squared error, in the same response units as the fitted data. Lower values indicate smaller typical residuals. \\
MAE & mae & Lower is better & Mean absolute error, in the same response units as the fitted data. Lower values indicate smaller average absolute residuals. \\
NRMSE (mean) & mnrmse & Lower is better & RMSE normalized by the response mean. Lower values indicate smaller error relative to the response magnitude. \\
NRMSE (std. dev.) & nrmse & Lower is better & RMSE normalized by the response standard deviation. Lower values indicate smaller error relative to response variability. \\
AIC & aic & Lower is better & Akaike Information Criterion. Information-criterion metrics balance how well a model fits the data with model complexity (the number of fitted parameters). Among comparable fitted models, a lower AIC value generally indicates a better model representation: the model explains the data well without unnecessary complexity. Absolute AIC values, including negative values, should not be interpreted as standalone quality grades. \\
AIC\textsubscript{c} & aicc & Lower is better & Corrected Akaike Information Criterion. AIC\textsubscript{c} applies a small-sample correction to AIC and is often preferred when the number of fitted observations (n) is limited relative to the number of fitted parameters (k), with n/k < 40 commonly used as a rule of thumb. Among comparable fitted models, a lower AIC\textsubscript{c} value generally indicates a better balance of fit quality and model complexity; AIC\textsubscript{c} is undefined when there are too few observations relative to the number of parameters. \\
BIC & bic & Lower is better & Bayesian Information Criterion. BIC balances how well a model fits the data with model complexity, usually applying a stronger penalty for additional parameters than AIC. Among comparable fitted models, a lower BIC value generally indicates a better model representation without unnecessary complexity. Absolute BIC values, including negative values, should not be interpreted as standalone quality grades. \\
\bottomrule
\end{longtable}

For additional background on common error metrics, scale-dependent accuracy measures, and adjusted \(R^2\), see references \cite{hodson2022,hyndman2006,miles2014}.

For AIC\textsubscript{c}, a commonly cited rule of thumb is to prefer AIC\textsubscript{c} over AIC when the ratio of fitted observations to estimated parameters is small, such as \(n/k < 40\), where \(n\) is the number of fitted observations and \(k\) is the number of estimated model parameters \cite{burnham2002}.

\section{Uncertainty and error bars}

The following descriptions summarize how fitted-parameter uncertainty, fitting-plot bands, and prediction error bars are calculated or interpreted for each approach.

\subsection{Approach 1: Direct Correlation}

\textbf{Fitting uncertainty}

Parameter estimates: fitted parameters are reported as estimate \(\pm\) standard error. Standard errors are derived from the nonlinear least-squares parameter covariance matrix. Correlations among fitted parameters are retained when uncertainty-aware parameter values are created.

Fitting plot bands: shaded 95\% prediction bands are generated by Monte Carlo sampling of the fitted parameter covariance matrix. The sampled model variance is combined with the residual variance from the fitted data, then converted to a two-sided interval using a t multiplier when residual degrees of freedom are available.

These bands describe approximate prediction uncertainty around the fitted relationship, not regulatory acceptance limits. They depend on the fitted model form, local covariance approximation, residual behavior, and preprocessing choices.

\textbf{Prediction uncertainty}

Prediction plot error bars: when a prediction dataset is supplied, predicted in vivo times are plotted with horizontal error bars. The bars are calculated from the same 95\% prediction-band calculation used for the fitted correlation model; the upper and lower bar lengths are the distances from the band mean to the corresponding upper and lower interval bounds at each input prediction time.

The error bars therefore include propagated fitted-parameter uncertainty and residual/model error from the Approach 1 correlation fit.

\subsection{Approach 2: Independent Model Fitting}

\textbf{Fitting uncertainty}

Parameter estimates: each in vitro and in vivo model parameter is reported as estimate \(\pm\) standard error. Standard errors are derived from the corresponding nonlinear least-squares covariance matrix, with fitted-parameter correlations retained in uncertainty-aware values.

Time constants: tau values, when available, are calculated from the fitted model and parameter covariance. The reported tau uncertainty is the propagated standard uncertainty from that covariance-based calculation.

Fitting plot bands: the in vitro and in vivo shaded 95\% prediction bands are generated separately. For each fitted model, the app Monte Carlo samples the fitted parameter covariance matrix, combines sampled model variance with residual variance, and converts the resulting standard uncertainty to a two-sided interval using a t multiplier when residual degrees of freedom are available.

\textbf{Prediction uncertainty}

Value-ratio prediction error bars: uncertainty is propagated from both fitted models. The app samples each model's parameter covariance matrix, calculates the in vivo/in vitro response ratio for each sample, combines ratio variance with residual-based ratio uncertainty, and converts that uncertainty to a two-sided 95\% interval. The plotted vertical error bars are scaled from the resulting ratio interval at each prediction time.

Time-constant-ratio prediction error bars: when both tau values are valid, uncertainty in each tau is propagated through the tau ratio using uncertainty-aware arithmetic. The standard uncertainty of the scaled prediction time is multiplied by a t value, where available, and plotted as a symmetric horizontal error bar.

Ratio-based uncertainty can become unstable when the denominator model prediction or denominator time constant is near zero. Predictions in those regions should be interpreted cautiously.

\subsection{Approach 3: Direct Mapping}

\textbf{Fitting uncertainty}

Parameter uncertainty: Approach 3 does not fit a parametric model, so there is no fitted-parameter covariance matrix, no parameter standard errors, and no model-based 95\% prediction band.

Fitting plots: the plots show interpolation-based ratios/mappings calculated from the preprocessed datasets. They should be interpreted as deterministic mappings from the supplied data and preprocessing choices, not as statistical uncertainty bands.

\textbf{Prediction uncertainty}

Prediction plot error bars: Approach 3 currently does not generate formal propagated error bars for prediction outputs because no parametric covariance model is fitted.

Uncertainty is therefore limited to the quality of the input data, preprocessing, interpolation, and the assumption that the observed mapping remains valid for the prediction dataset. Extrapolated points outside the interpolation range are excluded or returned as unavailable.

\section{Command-line interface}

The repository includes cli.py for batch-oriented fitting from the command line. Use python cli.py --help to see available arguments for input files, approaches, models, metrics, preprocessing options, cross-validation settings, and grid-search settings. The internal keys specified in this user guide are used as settings with cli.py.

\section{Extending the app}

The IVIVC app is designed so that common extensions can be added with minimal changes to the main application code. In many cases, new models, performance metrics, and preprocessing methods can be added by following the structure of the existing files in the relevant subdirectories.

